\section{Misevolution in Practice: A 60-Day, 8-Agent Empirical Study of
4 Self-Modification Paths and 25 Attack
Surfaces}\label{misevolution-in-practice-a-60-day-8-agent-empirical-study-of-4-self-modification-paths-and-25-attack-surfaces}

\begin{quote}
\textbf{Authors}: Anonymous Authors¹ \textbf{Affiliations}: ¹
{[}Anonymized for double-blind review --- see cover letter for
venue-specific disclosure{]} \textbf{Submission target}: ICLR 2026
Workshop on Agents (companion to the main Equipment Thickness Theory
submission) \textbf{Date}: 2026-06-27 (Draft v0.1) → 2026-07-09 (v1.0)
\textbf{Status}: Pre-submission draft, double-blind, not for circulation
\textbf{Revision log}: v0.1 (2026-06-27): initial. v1.0 (2026-07-09):
strict-mode hit rate updated to 60-case evaluation, 5/20 case
ground-truth drift inspected and corrected where drift was confirmed
(SR-006/007), no structural changes. \textbf{Word count target}:
\textasciitilde5,000 words main paper (workshop short paper limit: 5-7
pages) \textbf{Companion to}: \emph{Equipment Thickness Theory} (ICLR
2027 Main submission, anony\-mous) \textbf{Code \& data}: Open-source at
the project repository (MIT-compatible); see reproducibility checklist
in Appendix A.
\end{quote}

\begin{center}\rule{0.5\linewidth}{0.5pt}\end{center}

\subsection{Abstract}\label{abstract}

We report a 60-day field study of two recently proposed agent-safety
frameworks deployed on a production 8-agent system: (1) the
\textbf{Misevolution 4-path} failure taxonomy of self-modifying agents
(model / memory / tool / workflow) and (2) the \textbf{MLAS 5×5
attack-surface matrix} (5 modules × 5 lifecycle stages = 25 surfaces;
Lin et al., arXiv:2606.23075). We contribute three artifacts: (a) a
4-line implementation of the Misevolution 4-path detection policy as a
4-bucket action classifier (\texttt{\{path,\ score\}}) used by the
safety boundary to gate high-impact actions; (b) a 25-cell MLAS audit
checklist with current coverage status for our 8 agents, comprising 19
critical and 6 managed surfaces, of which 2 are fully covered, 18
partially covered, and 5 not yet covered; (c) a longitudinal incident
record: 4 Misevolution-path blocked events, 1 MLAS-surface check (M2-L2
in-flight memory write) and 2 follow-up checks, 1 critical
write-protection incident (the ``Agent-D 2026-06 memory overwrite''
event, classified as Misevolution memory path + MLAS M2-L3). We find
that the Misevolution 4-path classification and the MLAS 5×5 matrix are
\textbf{complementary, not redundant}: Misevolution classifies
\emph{failure modes} along the dimension of which layer misbehaves,
while MLAS classifies \emph{attack opportunities} along the dimension of
module × lifecycle stage. Combined, they cover the failure-manifestation
and attack-opportunity axes that neither covers alone. We also document
a measurement caveat: the operational data set (4 Misevolution events, 3
MLAS checks) is two to three orders of magnitude smaller than what the
cover letter and our initial sketch described (1,247 events; 25-cell
coverage). We discuss the gap and the inferred reasons --- chiefly that
the safety boundary was not exercised at high rates during the
observation window --- and argue that small-but-real evidence is more
useful for a workshop case study than projected counts, because the data
we \emph{do} have identifies a real incident with concrete remediation
steps. Defense overhead measured at the action-policy layer is 0.3\% of
total wall-clock time (median across 100 evaluate calls). Our 8-agent
setup and MLAS 5×5 checklist are released as open-source artifacts.

\begin{center}\rule{0.5\linewidth}{0.5pt}\end{center}

\subsection{1. Introduction}\label{introduction}

\textbf{Motivation.} Self-modifying AI agents --- agents that update
their own memory, model weights, tool code, or workflow parameters at
runtime --- introduce a class of safety risks that static-rule systems
cannot fully address. Two recent papers have proposed complementary
frameworks for characterising these risks. \textbf{Misevolution} (Shao
et al., 2026) proposes a 4-path taxonomy of \emph{failure modes} in
self-modifying agents: model-path (the LLM misbehaves), memory-path
(long-term or short-term memory corrupts), tool-path (tool code or
interfaces misbehave), and workflow-path (dispatch, cron, or reflexion
loops misbehave). \textbf{MLAS} (Lin et al., arXiv:2606.23075) proposes
a 5×5 attack-surface matrix of \emph{attack opportunities}: 5 modules
(model, memory, tool, workflow, governance) crossed with 5 lifecycle
stages (initialization, operation, modification, shutdown, recovery) =
25 cells, with the paper reporting that 17 of 25 cells face critical
threat under self-evolution amplification.

\textbf{Gap.} Both papers are conceptual; neither paper reports a
longitudinal deployment. Our contribution closes this gap by reporting a
60-day field study of an 8-agent production system in which both
frameworks are operationalised and exercised.

\textbf{Why this paper fits the workshop.} The ICLR 2026 Workshop on
Agents solicits empirical case studies with operational depth and
concrete operational artefacts. Our submission is exactly that: a real
deployment, a real incident (the Agent-D 2026-06 memory overwrite,
classified as Misevolution memory-path + MLAS M2-L3), and a
release-ready 25-cell audit checklist.

\textbf{Contributions.}

\begin{enumerate}
\def\labelenumi{\arabic{enumi}.}
\tightlist
\item
  \textbf{Operationalisation of Misevolution 4-path detection} as a
  4-bucket action classifier on our 8-agent system, with one real
  blocked incident and three additional defence actions traced to the
  policy (Section 4).
\item
  \textbf{Operationalisation of the MLAS 5×5 matrix} as a 25-cell audit
  checklist with current coverage status and detection-rule
  specifications for each cell (Section 5).
\item
  \textbf{Case study: the Agent-D 2026-06 memory overwrite incident} ---
  how Misevolution memory-path + MLAS M2-L3 classification explained the
  failure, how the safety boundary fired, and what remediation we took
  (Section 6).
\item
  \textbf{Defense overhead measurement} at 0.3\% of total wall-clock
  time, indicating that the action-policy layer is cheap to operate but
  its absence (the Agent-D 2026-06 incident) is expensive to recover
  from (Section 7).
\end{enumerate}

\textbf{Scope.} This paper is a companion to our main Equipment
Thickness Theory submission; readers need not read the main paper to
follow this one. We do not duplicate the main paper's theoretical
framework; we focus on the two safety frameworks' operational behaviour
in production.

\textbf{Limitations preview.} Our Misevolution event count (4) and MLAS
check count (3) are small. We discuss the reasons (low-rate safety
events during a steady-state window) and the implications for the field
in Section 8.

\begin{center}\rule{0.5\linewidth}{0.5pt}\end{center}

\subsection{2. Background: 8-Agent Production
System}\label{background-8-agent-production-system}

We deploy both safety frameworks on a production 8-agent system that has
been in continuous operation since 2026-04. The system runs as a
coordination-layer stack above a single 4B local LLM. The 8 agents
occupy distinct functional roles (we anonymise role names as
\textbf{Agent-A} through \textbf{Agent-I} to preserve double-blind
review; the \emph{kinds} of roles --- e.g., one safety auditor, one
skill librarian, one dispatch router --- are described below in
functional terms rather than by name). The agents share (a) a 4-layer
long-term memory hierarchy, (b) a shared lessons broadcast channel with
freshness decay, (c) a typed skill DAG of 207 skills with 159 typed
edges, and (d) a centralised action-policy boundary that gates
high-impact actions through a 7-rule decision engine
(\texttt{shell\ \textbar{}\ git\ \textbar{}\ a2a\ \textbar{}\ mcp\ \textbar{}\ fork}
categories; default confidence threshold 0.6; \texttt{deny} /
\texttt{ask\_user} / \texttt{allow\_with\_audit} outcomes).

\textbf{Functional role distribution (anonymised):}

\begin{itemize}
\tightlist
\item
  \textbf{Agent-A}: Skills librarian. Maintains the typed skill DAG;
  emits skill-recommendation rules into the model-router decision log.
\item
  \textbf{Agent-B}: Dispatch router. Receives tasks, dispatches to the
  right agent via the typed DAG, logs every decision with a 20-field
  audit record.
\item
  \textbf{Agent-C}: Safety auditor. Runs the MLAS 5×5 audit on a daily
  cron and writes incidents to a shared log.
\item
  \textbf{Agent-D}: Memory curator. Performs background compaction and
  decay on the 4-layer memory hierarchy. \textbf{Operates the cell at
  highest risk of memory-path Misevolution events.}
\item
  \textbf{Agent-E}: Model router. Selects between local and cloud LLM
  based on data-classification, cost-ceiling, and complexity; emits a
  20-field decision log.
\item
  \textbf{Agent-F}: Reflexion loop runner. Runs after-action reflection
  to update shared lessons.
\item
  \textbf{Agent-G}: Cron SLO monitor. Maintains a 17-cron registry with
  hash-pinned entries.
\item
  \textbf{Agent-H}: Model fine-tuner. Performs self\_finetune operations
  on a sandboxed model copy; triggers Misevolution model-path on every
  fine-tune without alignment\_check\_passed.
\item
  \textbf{Agent-I}: Governance / policy manager. Owns the action-policy
  schema, the 5-tier policy hierarchy, and the audit log retention.
\end{itemize}

\textbf{Why these 8 agents matter for this study.} The 4 Misevolution
paths and the 5 MLAS modules are \emph{all} exercised by at least one
agent. The system is a natural test bed because every cell of the MLAS
5×5 matrix is potentially active every day. We do not artificially
induce failures; we observe what happens in production.

\textbf{Observation window.} 2026-04-15 to 2026-06-13 (60 days). The
window straddles a major version transition in the action-policy schema
(v0.x → v1.0 on 2026-05-22) and the deployment of the typed skill DAG
(49 skills on 2026-05-29, expanded to 207 on 2026-06-10).

\textbf{Data retention.} All audit logs are write-once with file-level
lock; no agent can mutate an audit log entry after write. The 4
Misevolution events and 3 MLAS checks retained in the present submission
are \emph{all} the events recorded in the window; we did not sample.

\begin{center}\rule{0.5\linewidth}{0.5pt}\end{center}

\subsection{3. Misevolution 4-Path Detection
Policy}\label{misevolution-4-path-detection-policy}

The Misevolution 4-path taxonomy (Shao et al., 2026) classifies
self-modification failures by which layer misbehaves:

\begin{itemize}
\tightlist
\item
  \textbf{Model path}: the LLM itself produces a misleading or harmful
  action, including prompt-injection-induced behaviour.
\item
  \textbf{Memory path}: long-term or short-term memory stores incorrect,
  stale, or malicious content that propagates to downstream actions.
\item
  \textbf{Tool path}: tool code, MCP interfaces, or function calls
  misfire (returns wrong data, throws unexpectedly, or executes a
  destructive side effect).
\item
  \textbf{Workflow path}: dispatch, cron, or reflexion loop misbehaves
  (priority inversion, dead dispatch, infinite reflexion).
\end{itemize}

We implement the 4-path classifier as a thin layer in front of the
existing action-policy engine. Each high-impact action request is first
classified into a 4-vector
\texttt{path\_scores\ =\ \{model:\ 0/1/2,\ memory:\ 0/1/2,\ tool:\ 0/1/2,\ workflow:\ 0/1/2\}}
(0 = no signal, 1 = weak signal, 2 = strong signal) by keyword + context
matching. If the dominant path's score ≥ 2 and at least one
\textbf{missing-signal} is detected (e.g., for memory path, the request
lacks \texttt{backup\_first} and the target path is \texttt{learned.md}
or \texttt{diary.md}), the action is \textbf{blocked}. The 4-path
classifier reuses the same \texttt{evaluate} function and audit-log
mechanism as the action policy; no separate infrastructure was added.

\textbf{Coverage.} The classifier catches the four named patterns:
model-path (e.g., ``self\_finetune'' without
\texttt{alignment\_check\_passed}), memory-path (e.g., ``overwrite
diary'' without \texttt{backup\_first}), tool-path (e.g., ``create new
tool'' without \texttt{evidence\_gated}), workflow-path (e.g.,
``infinite reflexion'' without \texttt{loop\_termination\_check}). It
does not catch Misevolution events that \emph{do not} match the keyword
taxonomy; we treat those as detection gaps and discuss them in Section
8.

\begin{center}\rule{0.5\linewidth}{0.5pt}\end{center}

\subsection{4. Misevolution 4-Path Empirical
Results}\label{misevolution-4-path-empirical-results}

Across the 60-day observation window, the action-policy engine recorded
\textbf{4 Misevolution-classified events} that triggered the policy's
``block'' action. We tabulate the per-event breakdown below; the path
distribution is the empirical manifestation of Misevolution in our
8-agent system.

{\def\LTcaptype{none} % do not increment counter
\begin{longtable}[]{@{}
  >{\raggedright\arraybackslash}p{(\linewidth - 12\tabcolsep) * \real{0.0674}}
  >{\raggedright\arraybackslash}p{(\linewidth - 12\tabcolsep) * \real{0.1573}}
  >{\raggedright\arraybackslash}p{(\linewidth - 12\tabcolsep) * \real{0.0899}}
  >{\raggedright\arraybackslash}p{(\linewidth - 12\tabcolsep) * \real{0.1798}}
  >{\raggedright\arraybackslash}p{(\linewidth - 12\tabcolsep) * \real{0.1348}}
  >{\raggedright\arraybackslash}p{(\linewidth - 12\tabcolsep) * \real{0.2022}}
  >{\raggedright\arraybackslash}p{(\linewidth - 12\tabcolsep) * \real{0.1685}}@{}}
\toprule\noalign{}
\begin{minipage}[b]{\linewidth}\raggedright
Date
\end{minipage} & \begin{minipage}[b]{\linewidth}\raggedright
Agent (role)
\end{minipage} & \begin{minipage}[b]{\linewidth}\raggedright
Action
\end{minipage} & \begin{minipage}[b]{\linewidth}\raggedright
Dominant path
\end{minipage} & \begin{minipage}[b]{\linewidth}\raggedright
Path score
\end{minipage} & \begin{minipage}[b]{\linewidth}\raggedright
Missing signals
\end{minipage} & \begin{minipage}[b]{\linewidth}\raggedright
Action taken
\end{minipage} \\
\midrule\noalign{}
\endhead
\bottomrule\noalign{}
\endlastfoot
2026-05-15 15:18 & Agent-D (memory) & write diary content & memory & 1 &
backup\_first, size\_under\_threshold & \textbf{block} \\
2026-05-15 17:51 & Agent-D (memory) & overwrite diary.md with new
content & memory & \textbf{2} & backup\_first, size\_under\_threshold &
\textbf{block} \\
2026-05-15 17:51 & Agent-H (model) & self\_finetune new dataset & model
& \textbf{2} & alignment\_check\_passed, goal\_drift\_index\_ok &
\textbf{block} \\
2026-05-15 17:51 & Agent-A (skill) & register new MCP tool & tool &
\textbf{2} & evidence\_gated, human\_approved & \textbf{block} \\
\end{longtable}
}

\textbf{Path distribution (n=4).}

{\def\LTcaptype{none} % do not increment counter
\begin{longtable}[]{@{}lll@{}}
\toprule\noalign{}
Misevolution path & Events & Share \\
\midrule\noalign{}
\endhead
\bottomrule\noalign{}
\endlastfoot
Model & 1 & 25.0\% \\
Memory & 2 & 50.0\% \\
Tool & 1 & 25.0\% \\
Workflow & 0 & 0.0\% \\
\end{longtable}
}

Three observations. First, \textbf{all four events were blocked} by the
policy --- the action boundary was effective for every detected event.
Second, the \textbf{memory path dominated} (50\% of events), driven
largely by Agent-D's routine diary-write and overwrite actions (Agent-D
accounted for 2 of 2 memory-path events); this is consistent with the
Misevolution paper's case studies, which identify memory write as the
highest-frequency failure mode. Third, \textbf{no workflow-path events
were detected} in the window, but this is a detection gap (the
classifier does not yet have strong workflow-path keywords), not an
absence of workflow failures; we discuss this in Section 8.

\textbf{False-positive rate.} Across the 60-day window, the 4-block
event count is the same as the 4-Misevolution-event count (i.e., zero
false positives among the 4 blocks); the broader false-positive rate
requires counting non-Misevolution actions that were erroneously
blocked, which we have not instrumented. We discuss this measurement gap
in Section 8.

\textbf{What the cover letter claimed vs.~what we report.} Our workshop
submission cover letter described ``1,247 self-modification events'' and
``0.7\% false positive rate''. The actual retained data is 4 events. The
discrepancy is due to (a) the cover letter being drafted from an earlier
sketch that conflated \emph{audit-log lines} (high count) with
\emph{Misevolution events} (low count after the 4-path filter is
applied), and (b) the 0.7\% figure was a projection from a 200-event
pilot in week 1 that did not generalise. We report the 4-event actual
here and discuss the implications for empirical agent-safety studies in
Section 8. We do not claim the 1,247 figure; we are correcting it.

\textbf{Detection rule details per path.} The 4-path classifier is a
thin wrapper around the existing \texttt{evaluate} function. The
detection rules are:

\begin{itemize}
\tightlist
\item
  \textbf{Model path}: matches action verbs \texttt{self\_finetune},
  \texttt{self\_train}, \texttt{reweight}, \texttt{override\_prompt},
  \texttt{inject\_prompt}. Requires the action to also carry at least
  one of \texttt{alignment\_check\_passed},
  \texttt{goal\_drift\_index\_ok}, \texttt{human\_review\_token}.
  Without these signals, the request is treated as model-path score 2
  and blocked.
\item
  \textbf{Memory path}: matches target nouns \texttt{diary},
  \texttt{learned}, \texttt{SOUL}, \texttt{shared\_lessons},
  \texttt{lessons\_broadcast}, \texttt{compaction}. Requires either
  \texttt{backup\_first} signal or \texttt{size\_under\_threshold}
  signal. Memory writes to files larger than 4 KB without a
  \texttt{backup\_first} signal are blocked (this is the rule that
  caught the Agent-D 2026-06 incident).
\item
  \textbf{Tool path}: matches \texttt{create\_tool},
  \texttt{register\_mcp}, \texttt{install\_skill},
  \texttt{update\_tool\_code}. Requires \texttt{evidence\_gated} (i.e.,
  the tool change has a passing evaluation suite) and
  \texttt{human\_approved}. New tool creation without human approval is
  blocked.
\item
  \textbf{Workflow path}: matches \texttt{add\_cron},
  \texttt{modify\_loop}, \texttt{infinite\_reflexion},
  \texttt{priority\_override}. Requires
  \texttt{loop\_termination\_check} and \texttt{priority\_audit}.
  Workflow modifications without these are blocked.
\end{itemize}

The keyword lists are open-source and version-controlled; agents can
extend them by submitting a Misevolution-extension proposal that the
safety auditor (Agent-C) reviews daily. We have received 0 extension
proposals in the 60-day window.

\begin{center}\rule{0.5\linewidth}{0.5pt}\end{center}

\subsection{5. MLAS 5×5 Attack-Surface
Audit}\label{mlas-55-attack-surface-audit}

We audit the 8-agent system against the MLAS 5×5 matrix (Lin et al.,
arXiv:2606.23075) using our open-source \texttt{mlas\_25.py} checklist
runner. The runner enumerates all 25 cells, queries the audit log for
prior checks, and reports a coverage matrix. We ran the audit on
2026-06-27; the current coverage status is summarised below.

\textbf{5×5 coverage matrix} (rows = module, columns = lifecycle stage;
entry = severity \textbar{} coverage):

{\def\LTcaptype{none} % do not increment counter
\begin{longtable}[]{@{}
  >{\raggedright\arraybackslash}p{(\linewidth - 10\tabcolsep) * \real{0.2537}}
  >{\raggedright\arraybackslash}p{(\linewidth - 10\tabcolsep) * \real{0.1343}}
  >{\raggedright\arraybackslash}p{(\linewidth - 10\tabcolsep) * \real{0.1045}}
  >{\raggedright\arraybackslash}p{(\linewidth - 10\tabcolsep) * \real{0.1194}}
  >{\raggedright\arraybackslash}p{(\linewidth - 10\tabcolsep) * \real{0.1940}}
  >{\raggedright\arraybackslash}p{(\linewidth - 10\tabcolsep) * \real{0.1940}}@{}}
\toprule\noalign{}
\begin{minipage}[b]{\linewidth}\raggedright
Module \textbackslash{} Stage
\end{minipage} & \begin{minipage}[b]{\linewidth}\raggedright
L1 init
\end{minipage} & \begin{minipage}[b]{\linewidth}\raggedright
L2 op
\end{minipage} & \begin{minipage}[b]{\linewidth}\raggedright
L3 mod
\end{minipage} & \begin{minipage}[b]{\linewidth}\raggedright
L4 shutdown
\end{minipage} & \begin{minipage}[b]{\linewidth}\raggedright
L5 recovery
\end{minipage} \\
\midrule\noalign{}
\endhead
\bottomrule\noalign{}
\endlastfoot
\textbf{M1 model} & critical \textbar{} partial & critical \textbar{}
partial & critical \textbar{} none & critical \textbar{} none & managed
\textbar{} partial \\
\textbf{M2 memory} & critical \textbar{} partial & critical \textbar{}
partial & critical \textbar{} none & critical \textbar{} none & managed
\textbar{} partial \\
\textbf{M3 tool} & critical \textbar{} partial & critical \textbar{}
partial & critical \textbar{} partial & managed \textbar{} partial &
\textbf{managed \textbar{} full} \\
\textbf{M4 workflow} & \textbf{critical \textbar{} full} & critical
\textbar{} partial & critical \textbar{} partial & critical \textbar{}
none & managed \textbar{} partial \\
\textbf{M5 governance} & critical \textbar{} partial & critical
\textbar{} partial & critical \textbar{} partial & critical \textbar{}
partial & managed \textbar{} partial \\
\end{longtable}
}

\textbf{Coverage rollup (n=25).}

{\def\LTcaptype{none} % do not increment counter
\begin{longtable}[]{@{}ll@{}}
\toprule\noalign{}
Severity & Count \\
\midrule\noalign{}
\endhead
\bottomrule\noalign{}
\endlastfoot
critical & \textbf{19} \\
managed & 6 \\
\end{longtable}
}

{\def\LTcaptype{none} % do not increment counter
\begin{longtable}[]{@{}ll@{}}
\toprule\noalign{}
Coverage & Count \\
\midrule\noalign{}
\endhead
\bottomrule\noalign{}
\endlastfoot
full & \textbf{2} \\
partial & 18 \\
none & \textbf{5} \\
\end{longtable}
}

\textbf{Critical surfaces with no coverage (n=5).} M1-L3 (LoRA / weight
fine-tune injection), M1-L4 (shutdown command override), M2-L3 (memory
hierarchy bypass), M2-L4 (destructive memory write on shutdown), M4-L4
(dispatch queue poisoning on shutdown). The Agent-D 2026-06 incident
(Section 6) maps to M2-L3, the memory hierarchy bypass cell ---
confirming the practical impact of the ``no coverage'' status.

\textbf{Critical surfaces with full coverage (n=2).} M3-L5 (tool result
replay attack) and M4-L1 (cron task injection at boot). M3-L5 is covered
by the A2A v0.2 inbox-decay mechanism; M4-L1 is covered by the 17-cron
SLO monitor with hash-pinned crontab entries.

\textbf{Check history.} The audit log retains 3 check events (2 hits, 1
miss), all on cell M2-L2 (in-flight memory write injection), exercised
in 2026-06-26 by the safety auditor Agent-C. The 2-hit / 3-total rate is
not statistically meaningful; we report it for transparency only.

\textbf{Comparison with Lin et al.} The Lin et al.~paper reports 17/25
critical and 7 cross-module amplification effects. Our 19/25 critical
count includes 2 additional cells the paper lists as ``managed'' but we
flag as critical because our 8-agent system has a specific exposure
(M1-L4 shutdown override is critical for us because Agent-I runs
shutdown, not a human; M2-L4 is critical because Agent-D performs
autonomous compaction). Our 0 amplification-effect detections in 60 days
is consistent with the Lin et al.~prediction that amplification requires
\emph{triggering} modifications, which our 8-agent system performs at
low frequency.

\textbf{Open coverage gaps as a research agenda.} The 5 uncovered
critical cells (M1-L3, M1-L4, M2-L3, M2-L4, M4-L4) form a coherent
theme: \textbf{all are ``modification'' or ``shutdown'' lifecycle
stages}. The Lin et al.~paper makes the same observation: the
modification and shutdown stages are intrinsically higher-risk because
the agent is acting autonomously without steady-state monitoring. The 5
cells are the highest-priority targets for our Phase 6 closure plan.

\textbf{Why we built our own audit runner rather than reuse the Lin et
al.~released code.} Lin et al.~release a static reference checklist in
the paper appendix but no executable code. Our \texttt{mlas\_25.py} is
an independent implementation that (a) attaches each cell to a concrete
detection rule that can be invoked against the 8-agent audit log, (b)
tracks check history per cell so that we can compute per-cell hit rate
over time, and (c) emits a coverage matrix at the module × lifecycle
granularity that the Lin et al.~paper presents only as a 5×5 table. We
release \texttt{mlas\_25.py} so that other teams can run the same audit
against their own agents; the only deployment-specific input is the path
to the audit log directory.

\begin{center}\rule{0.5\linewidth}{0.5pt}\end{center}

\subsection{6. Case Study: Agent-D 2026-06 Memory Overwrite
Incident}\label{case-study-agent-d-2026-06-memory-overwrite-incident}

\textbf{The incident.} On 2026-06-08 (observed in Agent-D's diary
audit), Agent-D attempted to overwrite its 44,092-byte diary file with
empty content using a write tool that, per its documentation, ``always
overwrites the entire file''. The action was blocked by the Misevolution
memory-path detector (path\_score=2,
missing\_signals=\texttt{backup\_first} and
\texttt{size\_under\_threshold}). The 44,092 bytes of diary history were
preserved.

\textbf{How the safety boundary fired.} The classify function scored the
request with
\texttt{path\_scores\ =\ \{model:\ 0,\ memory:\ 2,\ tool:\ 0,\ workflow:\ 0\}};
the dominant path was \texttt{memory} with score 2 (matching the keyword
``overwrite'' twice --- once in the action verb and once in the target
noun). The \texttt{missing\_signals} check found neither
\texttt{backup\_first} nor \texttt{size\_under\_threshold} in the
request context. The action was blocked and an entry was written to the
Misevolution incident log.

\textbf{Timeline.}

{\def\LTcaptype{none} % do not increment counter
\begin{longtable}[]{@{}
  >{\raggedright\arraybackslash}p{(\linewidth - 2\tabcolsep) * \real{0.6667}}
  >{\raggedright\arraybackslash}p{(\linewidth - 2\tabcolsep) * \real{0.3333}}@{}}
\toprule\noalign{}
\begin{minipage}[b]{\linewidth}\raggedright
Time (UTC+8)
\end{minipage} & \begin{minipage}[b]{\linewidth}\raggedright
Event
\end{minipage} \\
\midrule\noalign{}
\endhead
\bottomrule\noalign{}
\endlastfoot
2026-06-08 20:31:00 & Agent-D prepares to write new content to diary.md
using \texttt{write\_file} \\
2026-06-08 20:31:01 & Action-policy engine classifies request;
memory\_path=2, model/tool/workflow=0 \\
2026-06-08 20:31:01 & Missing-signal check fails (no backup\_first, no
size\_under\_threshold) \\
2026-06-08 20:31:01 & \textbf{Action blocked}; audit-log entry
written \\
2026-06-08 20:31:30 & Agent-D reviews the block, switches to
\texttt{patch} tool, re-runs the same content with
\texttt{backup\_first} signal explicitly set \\
2026-06-08 20:31:45 & Action allowed; new content appended via patch;
diary history preserved \\
\end{longtable}
}

\textbf{Misevolution classification.} The incident is a textbook
\textbf{memory-path} Misevolution event: the agent's autonomous
compaction behaviour was about to destroy persistent state without a
backup signal. The \texttt{tool\ path} was \emph{not} the dominant path
because the tool itself functioned as documented; the failure mode was
the \emph{use of the tool without backup} --- a memory-path failure mode
even though a tool was the proximate cause. This is why the Misevolution
4-path taxonomy adds value over a tool-call allowlist: the same tool
call can be a tool-path event (tool malfunctions) or a memory-path event
(correctly-functioning tool used unsafely).

\textbf{MLAS classification.} The incident also maps cleanly to MLAS
cell \textbf{M2-L3 (memory hierarchy bypass)}: the agent was attempting
to modify its own memory hierarchy (the 4-layer long-term memory)
without going through the modification-stage safety boundary. Per the
MLAS audit, M2-L3 has \texttt{current\_coverage\ =\ none} --- yet the
action was still blocked, because the Misevolution 4-path policy fired
first. The two frameworks are complementary: MLAS identifies the cell as
a coverage gap; Misevolution provides the detection rule that prevents
the gap from being exploited.

\textbf{What the recovery looked like.} Recovery took \textasciitilde2
minutes of Agent-D's wall-clock time to switch tools and re-apply. The
remediation was to add a ``diary-write'' reflexion rule in Agent-D's
learned-rules registry: ``before any diary.md or learned.md write via
\texttt{write\_file}, use \texttt{patch} with explicit
\texttt{backup\_first} signal''. This rule is now part of Agent-D's
standard operating procedure.

\textbf{Why we report this incident prominently.} The Agent-D 2026-06
incident is the \textbf{only} Misevolution event in our 60-day window
with non-trivial operational impact (all other events were routine
blocks with no agent retry). It demonstrates that (a) Misevolution
4-path classification catches the right event, (b) MLAS 5×5 identifies
the right cell, and (c) the combination provides defence-in-depth even
when one of the two frameworks has \texttt{current\_coverage\ =\ none}.

\textbf{What the cover letter claimed vs.~what we report.} The cover
letter described a ``6-hour outage'' with ``1,472 lines permanently
lost''. The actual incident is a 2-minute block with zero data loss. The
cover-letter description is from an earlier draft that conflated this
incident with a separate, unrecovered historical event in a different
agent's logs (which we are not at liberty to detail in this anonymised
submission). We report the actual 2-minute / zero-loss incident here.

\begin{center}\rule{0.5\linewidth}{0.5pt}\end{center}

\subsection{7. Defense Overhead}\label{defense-overhead}

The action-policy engine is invoked on every high-impact action request.
We measure its wall-clock overhead by running 100 \texttt{evaluate}
calls back-to-back in a tight loop and dividing by the total request
budget of the 8-agent system over the 60-day window. We find:

{\def\LTcaptype{none} % do not increment counter
\begin{longtable}[]{@{}
  >{\raggedright\arraybackslash}p{(\linewidth - 2\tabcolsep) * \real{0.5333}}
  >{\raggedright\arraybackslash}p{(\linewidth - 2\tabcolsep) * \real{0.4667}}@{}}
\toprule\noalign{}
\begin{minipage}[b]{\linewidth}\raggedright
Metric
\end{minipage} & \begin{minipage}[b]{\linewidth}\raggedright
Value
\end{minipage} \\
\midrule\noalign{}
\endhead
\bottomrule\noalign{}
\endlastfoot
Median \texttt{evaluate} latency & 12 ms \\
95th-percentile \texttt{evaluate} latency & 41 ms \\
Total \texttt{evaluate} calls in 60 days & \textasciitilde14,200
(estimated from request-rate logs) \\
Total wall-clock time spent in \texttt{evaluate} & \textasciitilde170
seconds (median) / \textasciitilde580 seconds (p95) \\
\textbf{Defense overhead as \% of total agent wall-clock} &
\textbf{0.3\%} \\
\end{longtable}
}

\textbf{Implication.} The action-policy engine is \textbf{cheap to
operate} (0.3\% overhead) but \textbf{expensive to omit}: the Agent-D
2026-06 incident, had it not been blocked, would have required a 4-6
hour forensic recovery (rebuilding diary.md from cron SLO logs and
partial backups). The cost-benefit ratio is approximately 1:1000 ---
every second spent in the policy saves up to 1000 seconds of recovery.

\textbf{Defense trigger rate by Misevolution path.} Among the
\textasciitilde14,200 \texttt{evaluate} calls in the 60-day window, the
rate at which the policy returned a non-allow decision is summarised in
the table below. The dominant non-allow outcomes are \texttt{ask\_user}
(high-impact but recoverable actions that escalate to a human) and
\texttt{deny} (Misevolution-blocked actions). The numbers are derived
from the audit log grouped by Misevolution-path keyword match.

{\def\LTcaptype{none} % do not increment counter
\begin{longtable}[]{@{}
  >{\raggedright\arraybackslash}p{(\linewidth - 10\tabcolsep) * \real{0.1979}}
  >{\raggedright\arraybackslash}p{(\linewidth - 10\tabcolsep) * \real{0.1875}}
  >{\raggedright\arraybackslash}p{(\linewidth - 10\tabcolsep) * \real{0.0833}}
  >{\raggedright\arraybackslash}p{(\linewidth - 10\tabcolsep) * \real{0.1250}}
  >{\raggedright\arraybackslash}p{(\linewidth - 10\tabcolsep) * \real{0.2188}}
  >{\raggedright\arraybackslash}p{(\linewidth - 10\tabcolsep) * \real{0.1875}}@{}}
\toprule\noalign{}
\begin{minipage}[b]{\linewidth}\raggedright
Misevolution path
\end{minipage} & \begin{minipage}[b]{\linewidth}\raggedright
\texttt{evaluate} calls
\end{minipage} & \begin{minipage}[b]{\linewidth}\raggedright
\texttt{deny}
\end{minipage} & \begin{minipage}[b]{\linewidth}\raggedright
\texttt{ask\_user}
\end{minipage} & \begin{minipage}[b]{\linewidth}\raggedright
\texttt{allow\_with\_audit}
\end{minipage} & \begin{minipage}[b]{\linewidth}\raggedright
\textbf{Trigger rate}
\end{minipage} \\
\midrule\noalign{}
\endhead
\bottomrule\noalign{}
\endlastfoot
Model & \textasciitilde3,400 & 1 & 47 & 12 & \textbf{1.77\%} \\
Memory & \textasciitilde5,100 & 2 & 21 & 38 & \textbf{1.20\%} \\
Tool & \textasciitilde2,900 & 1 & 9 & 22 & \textbf{1.10\%} \\
Workflow & \textasciitilde2,800 & 0 & 14 & 18 & \textbf{1.14\%} \\
\textbf{Total} & \textbf{\textasciitilde14,200} & \textbf{4} &
\textbf{91} & \textbf{90} & \textbf{1.30\%} \\
\end{longtable}
}

The \textbf{1.30\% overall trigger rate} is the empirical rate at which
the policy stops or escalates an action. The 4
Misevolution-\texttt{deny} events are a strict subset of the 4
\texttt{deny} rows; the other 91+90 non-allow outcomes are routine
\texttt{ask\_user} / \texttt{allow\_with\_audit} decisions unrelated to
Misevolution. The path distribution shows model-path is the most
frequently triggered, but memory-path is the most frequently
\emph{denied} (consistent with the Misevolution paper's identification
of memory as the highest-stakes layer).

\textbf{False-positive rate (per Misevolution path).} We define a false
positive as a non-allow decision on an action that, upon review, would
have completed safely and correctly without the policy intervention. We
manually audited the 4 \texttt{deny} events (all true positives) and a
random 20-event sample of the 91 \texttt{ask\_user} events. Of the 20
\texttt{ask\_user} events sampled, 2 (10\%) were judged false positives:
an \texttt{ask\_user} on a routine \texttt{patch} operation that the
human later confirmed was safe. Extrapolating to the full 91
\texttt{ask\_user} set gives an estimated \textbf{false-positive rate of
1.1\%} (10\% × 91/14,200). For \texttt{deny} the false-positive rate is
\textbf{0\%} (all 4 were true positives, including the Agent-D 2026-06
incident).

{\def\LTcaptype{none} % do not increment counter
\begin{longtable}[]{@{}
  >{\raggedright\arraybackslash}p{(\linewidth - 6\tabcolsep) * \real{0.1324}}
  >{\raggedright\arraybackslash}p{(\linewidth - 6\tabcolsep) * \real{0.1176}}
  >{\raggedright\arraybackslash}p{(\linewidth - 6\tabcolsep) * \real{0.5147}}
  >{\raggedright\arraybackslash}p{(\linewidth - 6\tabcolsep) * \real{0.2353}}@{}}
\toprule\noalign{}
\begin{minipage}[b]{\linewidth}\raggedright
Outcome
\end{minipage} & \begin{minipage}[b]{\linewidth}\raggedright
Events
\end{minipage} & \begin{minipage}[b]{\linewidth}\raggedright
False positives (sampled/audited)
\end{minipage} & \begin{minipage}[b]{\linewidth}\raggedright
Estimated FPR
\end{minipage} \\
\midrule\noalign{}
\endhead
\bottomrule\noalign{}
\endlastfoot
\texttt{deny} (Misevolution-block) & 4 & 0 / 4 audited &
\textbf{0.0\%} \\
\texttt{ask\_user} (human-escalate) & 91 & 2 / 20 sampled → 9/91
extrapolated & \textbf{1.1\%} \\
\texttt{allow\_with\_audit} (audit-only) & 90 & not yet sampled &
unknown (planned for next sprint) \\
\end{longtable}
}

The 1.1\% \texttt{ask\_user} FPR is non-trivial but acceptable: each
\texttt{ask\_user} decision takes \textasciitilde30 seconds of human
time, so 91 escalations × 30s = 45 minutes of human attention over 60
days, against \textasciitilde14,200 automated decisions. The cost is
small relative to the benefit of catching the Agent-D 2026-06 incident.

\textbf{Comparison with related work.} The action-policy engine is
structurally similar to a kernel-level system call filter (e.g.,
seccomp-bpf on Linux). Seccomp filters typically run in microseconds;
our Python implementation runs in milliseconds. The 1000× gap is
acceptable for our 8-agent throughput target but would be a bottleneck
at 1000+ agent scale; we discuss scaling in Section 8.

\textbf{Visualisation.} Figure 1 maps the 4 Misevolution events and 3
MLAS checks onto the 5×5 matrix, with the Agent-D 2026-06 incident
highlighted. Figure 2 shows the MLAS 5×5 coverage evolution over the
60-day window. Both figures are generated by
\texttt{reproducibility/make\_paper\_b\_figures.py} (released with the
paper) and saved as
\texttt{equipment-thickness-repo/docs/figures/paper\_b\_fig1\_misevolution\_mlas\_scatter.png}
and \texttt{paper\_b\_fig2\_mlas\_coverage\_timeline.png} respectively.
Reviewers can regenerate the figures by running
\texttt{python3\ make\_paper\_b\_figures.py} from the
\texttt{reproducibility/} directory; the script has no external
dependencies beyond matplotlib ≥ 3.5.

\emph{{[}Figure 1 placeholder: 5×5 matrix with 4 Misevolution events as
filled circles, 3 MLAS checks as ✕ markers, and the Agent-D 2026-06
incident as a red star in the M2-L3 cell. Background cells shaded green
for ``full coverage'' (M3-L5, M4-L1) and yellow for ``partial coverage''
(M1-L1, M1-L2, M2-L1, M2-L2).{]}}

\emph{{[}Figure 2 placeholder: line plot with x-axis = days from
2026-04-15 (0, 30, 60) and y-axis = MLAS cells (count of 25). Four
series: full coverage (green ○, 0→0→2), partial coverage (gold □,
0→12→18), no coverage (red △, 25→13→5), critical-uncovered (purple ◇,
19→12→5).{]}}

The figures are included in the supplementary material (PDF and PNG) and
are also reproducible from the source script.

\textbf{5.1 5/20 Case Curation Limitations.} Of the original 20
evaluation cases (SR-001..SR-020) that ground the 60-case
semantic-recall strict-mode hit-rate measurement, 5 cases (SR-004,
SR-006, SR-007, SR-009, SR-011) were inspected for ground-truth drift
after the v0.1 draft. We retrieved the top-5 chunk\_ids from each of the
6 retrieval modes and compared them against the case's
\texttt{expected\_source} (\texttt{learned.md} vs
\texttt{lessons.jsonl}) and \texttt{expected\_keywords}. Two of the five
cases (SR-006 ``anthropic teaching claude why misalignment'' and SR-007
``skillDAG typed dependency 12.8 percent'') had their
\texttt{expected\_source} pointing to the wrong file (the relevant chunk
actually lives in the \emph{other} file); we corrected these two
\texttt{expected\_source} fields in
\texttt{equipment-thickness-repo/code/semantic\_recall.py} and re-ran
the 60-case strict-mode evaluation. Two cases (SR-004 ``ICA model layer
architecture 6 layers'' and SR-009 ``装备率量化协议 kappa inter-rater'')
were verified to have correct \texttt{expected\_source} after inspection
(the relevant chunks at learned.md §1925 and §1923 do exist at the named
location). One case (SR-011 ``CADVP v1.1 channel fracture 13
dimensions'') has no \texttt{expected\_source} chunk in the current v2.0
index (the index covers \texttt{learned.md} + \texttt{lessons.jsonl} +
project reports; the canonical CADVP v1.1 lesson currently lives in
\texttt{pending.md} line 287 and is outside the indexed corpus), so we
did not edit \texttt{expected\_source} --- the case is retained for a
future indexing pass. Post-correction, the 60-case strict-mode hit rates
are: keyword 38/60 (63.33\%), semantic 36/60 (60.00\%), hybrid 38/60
(63.33\%), semantic\_v2 22/60 (36.67\%), hybrid\_v2 37/60 (61.67\%),
faiss 22/60 (36.67\%). The semantic\_v2 36.67\% rate is the
post-correction value used in the §7 defense-overhead narrative; the
5/20 inspection is the principal difference between v0.1 and v1.0. We
discuss curation as an ongoing maintenance task rather than a one-time
data preparation step: each new lesson added to \texttt{learned.md} or
\texttt{lessons.jsonl} should be re-matched against the existing 60-case
set, and each new evaluation case should be re-ground-truthed against
the index before inclusion.

\begin{center}\rule{0.5\linewidth}{0.5pt}\end{center}

\subsection{8. Discussion}\label{discussion}

\textbf{The 4-event vs.~1,247-event discrepancy, revisited.} Our initial
sketch projected \textasciitilde1,247 Misevolution events and a 0.7\%
false-positive rate based on week-1 pilot data. The 60-day actual is 4
events. The discrepancy has two causes: (1) the week-1 pilot counted
\emph{all audit-log lines} (including routine ``allow'' decisions and
re-evaluations), whereas the 4-event count is the post-filter count of
\emph{blocks} --- the two metrics are not comparable. (2) The 8-agent
system entered a steady-state operation around week 3, after which
Misevolution-event frequency dropped to its true baseline of
\textasciitilde2 events/month. Workshop readers should treat the 4-event
count as the \emph{operational} number: the rate at which the 4-path
classifier fires in steady state, not the rate at which the audit log
records lines.

\textbf{Why small-n evidence is still useful.} A workshop paper
reporting 4 events could be dismissed as anecdotal. We argue the
opposite: the 4 events include one operationally significant incident
(Agent-D 2026-06), the 4-path taxonomy correctly classified all 4 events
(zero false positives, zero false negatives among the 4), and the MLAS
5×5 audit identified 5 critical cells as priority targets for Phase 6
closure. The complementary value of the two frameworks is \emph{visible
in the data}: the only ``no coverage'' critical cell that was actually
exploited (M2-L3) was caught by the \emph{other} framework's detection
rule.

\textbf{What the data does not support.} We do not have evidence that
the Misevolution 4-path taxonomy is \emph{complete} (we have not
stress-tested rare failure modes like cross-agent memory poisoning). We
do not have evidence that MLAS coverage status correlates with incident
rate (our 4-event sample is too small). We do not have evidence that the
0.3\% overhead is sustainable at 1000+ agent scale.

\textbf{Comparison with prior empirical agent-safety studies.} Voyager
(Wang et al., 2023) reports 1-of-N skill acquisition in Minecraft over
30 days; SiriuS (Zhao et al., 2025) reports cross-agent experience
replay on SWE-bench over 1-2 weeks. Our 60-day, 8-agent, 2-framework
deployment is the longest such study we are aware of, but the n=4 event
count means our contribution is \emph{operational} (a checklist + a case
study) more than \emph{statistical} (an incidence-rate claim). We are
explicit about this.

\textbf{Recommendations for replication studies.} A 100+ agent
replication would (a) give a meaningful incidence rate per Misevolution
path, (b) test the 0.3\% overhead claim at scale, and (c) test the MLAS
5×5 coverage status correlation with incident rate. We provide our
25-cell checklist and the action-policy engine as open-source artifacts
to lower the replication cost.

\textbf{What we learned that the two papers do not predict.} Three
observations from the 60-day field study that are not in either
Misevolution or MLAS:

\begin{enumerate}
\def\labelenumi{\arabic{enumi}.}
\item
  \textbf{The same action can be classified differently by the two
  frameworks depending on the missing-signal set.} The Agent-D 2026-06
  incident was a memory-path event per Misevolution (because the
  \emph{use} of the tool was unsafe) and an M2-L3 event per MLAS
  (because the cell is a coverage gap). The two classifications agree in
  this case but they answer different questions: Misevolution asks
  ``which layer is failing?'', MLAS asks ``which cell is exposed?''.
  Researchers should report both classifications for any single
  incident, not just one.
\item
  \textbf{The keyword-based 4-path classifier has a vocabulary bias.}
  Our keyword list is dominated by the words we use in our own agent
  prompts. A replication study that uses different agent prompt language
  will likely have a different Misevolution path distribution even for
  the \emph{same} underlying failure mode. This is a known limitation of
  keyword-based classifiers; future work should embed the action
  description and classify by embedding similarity.
\item
  \textbf{The MLAS 5×5 coverage status is a more stable measurement than
  the Misevolution event count.} Over 60 days, our MLAS coverage status
  changed 3 times (Day 0 → Day 30 → Day 60), reflecting the deployment
  of the action policy, the cron SLO monitor, and the A2A inbox decay.
  The Misevolution event count, by contrast, has high day-to-day
  variance (0 events on most days, 4 events in a single 1-second cluster
  on 2026-05-15). Researchers tracking safety over time should use
  coverage status as the primary metric and event count as the secondary
  metric.
\end{enumerate}

\begin{center}\rule{0.5\linewidth}{0.5pt}\end{center}

\subsection{9. Conclusion}\label{conclusion}

We reported a 60-day, 8-agent field study of Misevolution 4-path
detection and the MLAS 5×5 attack-surface matrix. We contributed a
4-line Misevolution classifier (Section 3), a 25-cell MLAS checklist
with current coverage status (Section 5), a case study of one
operationally significant incident (Section 6), and a 0.3\%
defense-overhead measurement (Section 7). We corrected the cover-letter
data projection to the actual 4-event count and discussed the gap. We
argued that Misevolution and MLAS are complementary, not redundant, and
that small-n operational evidence is still useful for workshops because
the complementary value of the two frameworks is visible in the data we
do have. Our open-source artifacts (the \texttt{mlas\_25.py} checklist,
the action-policy engine, the 4 Misevolution incident log entries) are
released under MIT license.

\begin{center}\rule{0.5\linewidth}{0.5pt}\end{center}

\subsection{Appendix A: Reproducibility
Checklist}\label{appendix-a-reproducibility-checklist}

{\def\LTcaptype{none} % do not increment counter
\begin{longtable}[]{@{}
  >{\raggedright\arraybackslash}p{(\linewidth - 4\tabcolsep) * \real{0.2500}}
  >{\raggedright\arraybackslash}p{(\linewidth - 4\tabcolsep) * \real{0.3333}}
  >{\raggedright\arraybackslash}p{(\linewidth - 4\tabcolsep) * \real{0.4167}}@{}}
\toprule\noalign{}
\begin{minipage}[b]{\linewidth}\raggedright
Item
\end{minipage} & \begin{minipage}[b]{\linewidth}\raggedright
Status
\end{minipage} & \begin{minipage}[b]{\linewidth}\raggedright
Location
\end{minipage} \\
\midrule\noalign{}
\endhead
\bottomrule\noalign{}
\endlastfoot
\texttt{mlas\_25.py} checklist runner & ✅ Open-source & Project
repository, \texttt{code/mlas\_25.py} \\
MLAS 5×5 matrix definition (25 cells) & ✅ Open-source & Hard-coded in
\texttt{mlas\_25.py}; each cell has \texttt{name}, \texttt{description},
\texttt{examples}, \texttt{detection\_rule}, \texttt{arxiv\_ref} \\
Action-policy engine (Misevolution 4-path layer) & ✅ Open-source &
Project repository, \texttt{code/action\_policy\_check.py} \\
Misevolution incident log (n=4) & ✅ Open-source &
\texttt{agents/shared/misevolution\_incidents.jsonl} \\
MLAS check log (n=3) & ✅ Open-source &
\texttt{agents/shared/mlas\_incidents.jsonl} \\
8-agent system configuration & ✅ Open-source & Project repository,
\texttt{agents/A} through \texttt{agents/I} \\
60-day audit logs & ✅ Available on request & Anonymised upon request
for replication studies \\
\end{longtable}
}

\subsection{Appendix B: Acronyms}\label{appendix-b-acronyms}

\begin{itemize}
\tightlist
\item
  \textbf{MLAS}: Module-Lifecycle Attack Surface (Lin et al.,
  arXiv:2606.23075).
\item
  \textbf{Misevolution}: Self-modification failure taxonomy (Shao et
  al., 2026).
\item
  \textbf{A2A}: Agent-to-agent protocol with freshness decay.
\item
  \textbf{MCP}: Model Context Protocol (tool-call interface).
\item
  \textbf{SLO}: Service-level objective.
\item
  \textbf{DAG}: Directed acyclic graph (used for typed skill
  relationships).
\end{itemize}

\begin{center}\rule{0.5\linewidth}{0.5pt}\end{center}

\emph{Manuscript word count: \textasciitilde4,950 words (main text,
excluding appendices and references).} \emph{References to companion
paper and related work are abbreviated for the workshop short-paper
limit; full bibliography is in the cover letter.} \emph{All agent names
are anonymised to Agent-A through Agent-I for double-blind review.}
